% =============================================================================
\section{Introduction}\label{sec:introduction}
% =============================================================================

% Figure 1 declared at the top of the Introduction so the page-width float
% lands at the top of page 2, rather than floating to page 3 and leaving
% flushbottom to stretch the Introduction paragraphs.
\begin{figure*}[!t]
\centering
\includegraphics[width=\textwidth]{figure1.png}
\caption{Pipeline for NHS critical-care surge planning, from
regional data and physics-informed forecasting to scenario-based robust
allocation and evaluation.}
\label{fig:pipeline}
\end{figure*}

NHS critical-care surge planning during epidemic waves balances the
cost of pre-positioning capacity against the risk of unmet demand while
keeping inter-regional transfers manageable across the seven NHS England
regions. COVID-19 placed sustained pressure on adult
mechanical-ventilation capacity, making this a coupled forecasting and
allocation problem under regime shift. Existing forecasting work spans
statistical and recurrent baselines~\cite{ajao2024deep}, graph-based
epidemic and hospital-demand models~\cite{gao2021stan,gao2023hoist},
and physics-informed neural networks~\cite{raissi2019physics,ajao2025hybrid,qian2025physics};
interval and quantile evaluation is also standard in epidemic forecast
hubs~\cite{bracher2021evaluating}. The closest ICU-occupancy comparator,
a bidirectional LSTM with lagged mobility features and MC-dropout
intervals on Lazio data~\cite{valente2022forecasting}, like most such
work targets a single health system rather than multi-region planning.
These forecasters are usually
optimised for symmetric statistical accuracy, although critical-care
planning has an asymmetric error cost: under-prediction risks stranded
patients, whereas over-prediction wastes prepared capacity.

On the allocation side, robust facility-location and health-resource
models handle demand uncertainty with MILP, two-stage optimisation, and
metaheuristics~\cite{eriskin2022robust,wan2023two,shamseddin2025optimization,deb2002fast}.
Bertsimas et al.~\cite{bertsimas2022where} couple DELPHI dynamics to a
bilinear MIP for vaccination-site location, while Luo and
Stellato~\cite{luo2024equitable} embed neural-ODE compartmental dynamics
inside a McCormick-relaxed facility-location model. Our setting differs
from Smart Predict-then-Optimize~\cite{elmachtoub2022smart}. The
forecaster is not trained through allocation regret, but its
multi-quantile pinball objective is cost-asymmetric and feeds a robust
optimisation layer. To our knowledge, no prior work combines
physics-informed forecasting of NHS regional critical-care demand with
cost-asymmetric quantile training and an exact robust linear-programming
allocation across the seven NHS England regions.

This paper presents a forecast-to-decision pipeline
(\Cref{fig:pipeline}) in four phases. \emph{Phase~A} pre-trains a
per-region PINN on the $SEI_aI_sHCRD$ ODE system; its twelve outputs
(eight compartmental states, four transmission and clinical parameters)
are frozen and passed to a temporal GRU head that emits the horizon-wise
quantiles $q^{(0.1)}$, $q^{(0.5)}$, $q^{(0.9)}$ under a multi-quantile
pinball loss whose $q^{(0.9)}$ branch penalises under-prediction
nine-to-one, with a level-and-trend anchor. \emph{Phase~B}
discretises these quantiles into three scenarios with weights
$\pi=(0.2,0.6,0.2)$. \emph{Phase~C} solves the allocation exactly as a
linear program (deterministic and tail-risk-robust) via PuLP$+$CBC
against three heuristic baselines on a shared LP-slave oracle.
\emph{Phase~D} reports operational metrics and a sensitivity analysis
over the forecaster, budget $B$, tail-risk weight $\lambda_3$, travel
cap $\tau$, and transfer cap $\kappa$.

The main contributions are as follows:
(1)~\textbf{Cost-asymmetric forecaster.} A multi-quantile pinball loss,
a level-and-trend anchor, and a leakage-free per-origin recalibration
reduce $h{=}14$ RMSE by $25\%$ relative to refitted ARIMA and $42\%$
relative to a per-region GRU, and make the model competitive across all
horizons (best at $h{=}7$--$21$ and level with refitted ARIMA at
$h{=}28$, where seasonal-naive leads).
(2)~\textbf{Robust LP allocation.} With unconstrained transfers the
hybrid-regional robust LP ties every budget-exhausting policy on
coverage and wins on transfer burden ($4.6\times$ below
population-proportional); constraining transfers (sparser network or
capped transfer service) breaks the tie, the proportional rules
concede up to $60\%$ more expected unmet demand, and the robust LP is
never beaten on coverage.
